% GENERATED FILE. Edit canonical lessons, then run npm run generate:books.
% canonical-source-hash: fnv1a64:a23fd707d480f495
% canonical-lessons: LA-C355-condio, LA-C355-condimentum, LA-C355-aspergo, LA-C355-infundo, LA-C355-colo

\chapter{Condīmenta}
\label{ch:la-condimenta}
\hlchaptermodality{355}

\begin{chapteropening}
\textbf{By the end of this chapter:} \emph{I can tell how a dish was seasoned: to season, a seasoning, to sprinkle, to pour in and to strain.}

Complete the last lesson of chapter 355: I can tell how a dish was seasoned: to season, a seasoning, to sprinkle, to pour in and to strain.
\end{chapteropening}

\section[condiō, condīre]{condiō, condīre — to season, to spice; to pickle}

\label{lesson:LA-C355-condio}



\begin{quote}
Before the new one: say the Latin for to fry, then the Latin for to roast.
\end{quote}

\subsection*{You'll want to know: condiō, condīre}
\textbf{condiō, condīre} — \textquotedblleft{}to season, to spice; to pickle\textquotedblright{}.

\begin{quote}
Coquus carnem pipere et sale condīvit.
\end{quote}

\begin{quote}
\emph{The cook seasoned the meat with pepper and salt.}
\end{quote}

\emph{condiō, condīre, condīvī, condītum}. First to preserve fruit in vinegar or salt, then to flavour.

\begin{cousinweb}
\textbf{condiō} is taken by Lewis \& Short as a side form of \emph{condō}, \textquotedblleft{}to put away, to store\textquotedblright{}: first to put fruit away in pickle, then to season.
\end{cousinweb}

\begin{grammarlens}[title={What you've built}]
One more word for telling, explaining and describing at B1.
\end{grammarlens}

\subsection*{Your turn}
\begin{itemize}
\raggedright
  \item \emph{Say it:} \emph{condiō, condīre}
  \item \emph{Say it:} \emph{condiō, condīre}, once more
  \item \emph{Say it:} a sentence of your own with \emph{condiō, condīre}
  \item \emph{Recall:} say the Latin for to fry, then the Latin for to roast, then say \emph{condiō, condīre} again
\end{itemize}

\begin{tcolorbox}[breakable,colback=teal!4,colframe=teal!35!black,title={Before you move on}]
What does condiō, condīre mean? (\textquotedblleft{}To season, to spice; to pickle\textquotedblright{}.) Say it once more.
\end{tcolorbox}
\section[condīmentum, condīmentī]{condīmentum, condīmentī — a seasoning, a spice}

\label{lesson:LA-C355-condimentum}



\begin{quote}
Before the new one: say the Latin for to roast, then the Latin for to season.
\end{quote}

\subsection*{You'll want to know: condīmentum, condīmentī}
\textbf{condīmentum, condīmentī} — \textquotedblleft{}a seasoning, a spice\textquotedblright{}.

\begin{quote}
Famēs optimum condīmentum est.
\end{quote}

\begin{quote}
\emph{Hunger is the best seasoning.}
\end{quote}

Neuter, from \emph{condiō}. Cicero passes on the saying that hunger seasons food. English \emph{condiment}.

\begin{cousinweb}
\textbf{condīmentum} is built on \emph{condiō}, \textquotedblleft{}to season\textquotedblright{} (EtymDB), with the ending \emph{-mentum}. English \emph{condiment} comes from it.
\end{cousinweb}

\begin{grammarlens}[title={What you've built}]
One more word for telling, explaining and describing at B1.
\end{grammarlens}

\subsection*{Your turn}
\begin{itemize}
\raggedright
  \item \emph{Say it:} \emph{condīmentum, condīmentī}
  \item \emph{Say it:} \emph{condīmentum, condīmentī}, once more
  \item \emph{Say it:} a sentence of your own with \emph{condīmentum, condīmentī}
  \item \emph{Recall:} say the Latin for to roast, then the Latin for to season, then say \emph{condīmentum, condīmentī} again
\end{itemize}

\begin{tcolorbox}[breakable,colback=teal!4,colframe=teal!35!black,title={Before you move on}]
What does condīmentum, condīmentī mean? (\textquotedblleft{}A seasoning, a spice\textquotedblright{}.) Say it once more.
\end{tcolorbox}
\section[aspergō, aspergere]{aspergō, aspergere — to sprinkle}

\label{lesson:LA-C355-aspergo}



\begin{quote}
Before the new one: say the Latin for to season, then the Latin for a seasoning.
\end{quote}

\subsection*{You'll want to know: aspergō, aspergere}
\textbf{aspergō, aspergere} — \textquotedblleft{}to sprinkle\textquotedblright{}.

\begin{quote}
Coquus piper in carnem aspersit.
\end{quote}

\begin{quote}
\emph{The cook sprinkled pepper on the meat.}
\end{quote}

\emph{aspergō, aspergere, aspersī, aspersum}: \emph{ad} + \emph{spargō}, 'to scatter'. English \emph{asperse}.

\begin{cousinweb}
\textbf{aspergō} is built on \emph{spargō}, \textquotedblleft{}to scatter\textquotedblright{}, with \emph{ad-} (Lewis \& Short); the vowel changes inside the compound.
\end{cousinweb}

\begin{grammarlens}[title={What you've built}]
One more word for telling, explaining and describing at B1.
\end{grammarlens}

\subsection*{Your turn}
\begin{itemize}
\raggedright
  \item \emph{Say it:} \emph{aspergō, aspergere}
  \item \emph{Say it:} \emph{aspergō, aspergere}, once more
  \item \emph{Say it:} a sentence of your own with \emph{aspergō, aspergere}
  \item \emph{Recall:} say the Latin for to season, then the Latin for a seasoning, then say \emph{aspergō, aspergere} again
\end{itemize}

\begin{tcolorbox}[breakable,colback=teal!4,colframe=teal!35!black,title={Before you move on}]
What does aspergō, aspergere mean? (\textquotedblleft{}To sprinkle\textquotedblright{}.) Say it once more.
\end{tcolorbox}
\section[īnfundō, īnfundere]{īnfundō, īnfundere — to pour in, to pour on}

\label{lesson:LA-C355-infundo}



\begin{quote}
Before the new one: say the Latin for a seasoning, then the Latin for to sprinkle.
\end{quote}

\subsection*{You'll want to know: īnfundō, īnfundere}
\textbf{īnfundō, īnfundere} — \textquotedblleft{}to pour in, to pour on\textquotedblright{}.

\begin{quote}
Garum et oleum in patinam īnfūdit.
\end{quote}

\begin{quote}
\emph{He poured fish sauce and oil into the pan.}
\end{quote}

\emph{īnfundō, īnfundere, īnfūdī, īnfūsum}: \emph{in} + \emph{fundō}, 'to pour'. English \emph{infuse}.

\begin{cousinweb}
\textbf{īnfundō} is built on \emph{fundō}, \textquotedblleft{}to pour\textquotedblright{}, with \emph{in}, \textquotedblleft{}into\textquotedblright{} (Lewis \& Short). English \emph{infuse} comes from its participle.
\end{cousinweb}

\begin{grammarlens}[title={What you've built}]
One more word for telling, explaining and describing at B1.
\end{grammarlens}

\subsection*{Your turn}
\begin{itemize}
\raggedright
  \item \emph{Say it:} \emph{īnfundō, īnfundere}
  \item \emph{Say it:} \emph{īnfundō, īnfundere}, once more
  \item \emph{Say it:} a sentence of your own with \emph{īnfundō, īnfundere}
  \item \emph{Recall:} say the Latin for a seasoning, then the Latin for to sprinkle, then say \emph{īnfundō, īnfundere} again
\end{itemize}

\begin{tcolorbox}[breakable,colback=teal!4,colframe=teal!35!black,title={Before you move on}]
What does īnfundō, īnfundere mean? (\textquotedblleft{}To pour in, to pour on\textquotedblright{}.) Say it once more.
\end{tcolorbox}
\section[cōlō, cōlāre]{cōlō, cōlāre — to strain, to filter}

\label{lesson:LA-C355-colo}



\begin{quote}
Before the new one: say the Latin for to sprinkle, then the Latin for to pour in.
\end{quote}

\subsection*{You'll want to know: cōlō, cōlāre}
\textbf{cōlō, cōlāre} — \textquotedblleft{}to strain, to filter\textquotedblright{}.

\begin{quote}
Vīnum per linteum cōlāvit.
\end{quote}

\begin{quote}
\emph{He strained the wine through a cloth.}
\end{quote}

\emph{cōlō, cōlāre, cōlāvī, cōlātum}, from \emph{cōlum}, 'a strainer'. Not the verb \emph{colō, colere}, 'to farm'. English \emph{colander}.

\begin{cousinweb}
\textbf{cōlō} is built on \emph{cōlum}, \textquotedblleft{}a strainer, a colander\textquotedblright{} (EtymDB, Lewis \& Short). English \emph{colander} comes from its family.
\end{cousinweb}

\begin{grammarlens}[title={What you've built}]
One more word for telling, explaining and describing at B1.
\end{grammarlens}

\subsection*{Your turn}
\begin{itemize}
\raggedright
  \item \emph{Say it:} \emph{cōlō, cōlāre}
  \item \emph{Say it:} \emph{cōlō, cōlāre}, once more
  \item \emph{Say it:} a sentence of your own with \emph{cōlō, cōlāre}
  \item \emph{Recall:} say the Latin for to sprinkle, then the Latin for to pour in, then say \emph{cōlō, cōlāre} again
\end{itemize}

\begin{tcolorbox}[breakable,colback=teal!4,colframe=teal!35!black,title={Before you move on}]
What does cōlō, cōlāre mean? (\textquotedblleft{}To strain, to filter\textquotedblright{}.) Say it once more.
\end{tcolorbox}
